% Осень26!A30; нумерация внутри столбца без пустых строк.
% Источник: лекция 5, §8 (строка 238)
\ticket{27}{Измеримость функции непрерывной на множестве полной меры}
\label{ticket:27}

\subsection*{Измеримость функции, непрерывной на множестве полной меры}

Здесь лектор временно «переезжает» из общих пространств с мерой в \(\mathbb R^m\) с мерой Лебега \(\lambda\).

\begin{theorem}

Пусть \(E\subset\mathbb R^m\) измеримо по Лебегу, \(e\subset E\) и \(\lambda(e)=0\), \(E'=E\setminus e\). Пусть \(f:E\to\mathbb R\) (можно и в \(\overline{\mathbb R}\) --- лектор колебался, «чёрточка, не чёрточка, всё прекрасно работает») и \(f\) непрерывна \textbf{на \(E'\)} (в каждой точке \(E'\), относительно индуцированной топологии). Тогда \(f\) измерима \textbf{на \(E\)} (по Лебегу).

\end{theorem}

\begin{proof}

Проверим измеримость по множествам Лебега. В отличие от §2, здесь годится \textbf{любое} стандартное множество; возьмём \(E(f<a)\), \(a\in\mathbb R\). Разобьём поиск точек по двум частям: \[E(f<a)=\underbrace{E'(f<a)}_{\text{ищем в }E'}\cup\underbrace{\bigl(e\cap E(f<a)\bigr)}_{\text{ищем в }e}.\]

\textbf{(i)} \(E'(f<a)=\{x\in E': f(x)<a\}\) --- прообраз открытого луча \((-\infty,a)\) при непрерывном на \(E'\) отображении \(f|_{E'}:E'\to\mathbb R\); по критерию непрерывности он \textbf{открыт в \(E'\)}. («Лягушки всюду зелёные»: открыто --- значит открыто \textbf{в объёмлющем множестве} \(E'\).) Множество, открытое в \(E'\), имеет вид \(O\cap E'\) с открытым \(O\subset\mathbb R^m\), а \(E'=E\setminus e\) измеримо (\(E\) измеримо, \(e\) --- меры нуль, значит измеримо); пересечение измеримых измеримо. Итак, \(E'(f<a)\in\mathcal M^m\).

\textbf{(ii)} \(e\cap E(f<a)\subset e\), \(\lambda(e)=0\). Здесь выручает \textbf{полнота меры Лебега}: любое подмножество множества меры нуль измеримо (и имеет меру нуль).

Объединение двух измеримых множеств измеримо.

\end{proof}

\subsubsection*{Пример 1 : монотонные функции}

Пусть \(f:\mathbb R\to\mathbb R\) монотонна. По теореме 1-го семестра множество точек разрыва монотонной функции \textbf{не более чем счётно}: \(\exists\) (конечное или счётное) \(\{x_1,x_2,\dots\}\) такое, что \(f\) непрерывна на \(\mathbb R\setminus\{x_k\}\). Множество \(\{x_k\}\) имеет меру нуль, и по теореме \(f\) измерима.

Замечание лектора: «это скорее шутка» --- и так очевидно, что у монотонной функции множества Лебега устроены тривиально (лучи/промежутки); но теорема «заработала».

\subsubsection*{Пример 2 («настройка оптики»): характеристическая функция иррациональных чисел}

Пусть \(f=\chi_{\mathbb R\setminus\mathbb Q}\): \(f(x)=0\) при \(x\in\mathbb Q\), \(f(x)=1\) при \(x\notin\mathbb Q\). Функция непрерывна на \(\mathbb R\setminus\mathbb Q\) --- там она просто \textbf{постоянна} (\(\equiv1\)), а константа всяко непрерывна: \(|f(x)-f(a)|=0<\varepsilon\) безотносительно к «патологичности» самого множества. Мера \(\mathbb Q\) равна нулю, значит по теореме \(f\) измерима. (Впрочем, это и так ясно: \(\mathbb Q\) и \(\mathbb R\setminus\mathbb Q\) измеримы.)

\textbf{Мораль:} но на всём \(\mathbb R\) функция \textbf{всюду разрывна}: в любой окрестности любой точки она принимает и значение \(0\), и \(1\). То есть «непрерывность на почти всём множестве» \textbf{не} означает, что функция хороша на всём пространстве: на всём пространстве она «отличительная», а после сужения вдруг стала «хорошенькой». Этот эффект и нужен: множества Лебега разлагаются на части, среди которых кусочки меры нуль можно не контролировать.

\subsubsection*{Замечание : неполные меры и переопределение на множестве меры нуль}

Что будет, если отказаться от меры Лебега (и её полноты)? Пусть \((X,\mathcal C,\mu)\) --- общее пространство с мерой, \(e\subset X\), \(\mu(e)=0\), \(E'=X\setminus e\), и известно лишь, что \(f\) измерима \textbf{на \(E'\)}. Хотелось бы заключить «\(f\) измерима на \(X\)» --- \textbf{не получится}: в разложении \[X(f<a)=E'(f<a)\cup\bigl(e\cap X(f<a)\bigr)\] первое слагаемое измеримо, а второе --- \textbf{произвольное подмножество множества \(e\)}, выделенное неравенством; если \(f\) на \(e\) устроена патологически, оно может оказаться \textbf{неизмеримым} (меры, пришедшие из теоремы о лебеговском продолжении, полны, но, например, борелевская мера --- \textbf{нет}: у неё «упрощённая» область определения, с ней удобно работать, но нет полноты).

\textbf{Выход --- переопределение.} Функцию \(f\) можно \textbf{переопределить на \(e\)} так, чтобы получилась измеримая на всём \(X\) функция. Именно, положим (любая константа; лектор выбрал «весёлую») \[\tilde f(x):=2026\quad\text{при }x\in e,\qquad \tilde f|_{E'}:=f|_{E'}.\] Тогда \(X(\tilde f<a)=E'(f<a)\cup A_a\), где \(A_a=\varnothing\), если \(a\le2026\), и \(A_a=e\), если \(a>2026\) --- «добавка» либо пуста, либо совпадает с \(e\) целиком, т.е. всегда измерима (как множество меры нуль). Починили. \(\blacksquare\)
